\section{Critical target mass and posterior circuits}
\label{ac-appendix}

This appendix proves the finite-width critical input used in the main
theorem. Throughout, constants may depend on one fixed activation $\phi$.
We assume that $\phi$ is odd, strictly increasing, bounded by one,
and $1$-Lipschitz; that it is $C^5$ with bounded derivatives through
order five, concave on $(0,\infty)$, and satisfies
\[
 \phi'(0)=1,\qquad \phi'''(0)<0,\qquad
 \phi'(u)\ge c_0\operatorname{sech}^2u.
\]
Writing $\beta=\lim_{u\to\infty}\phi(u)>0$, assume
$\beta-\phi(u)\le C(1+u)^{-a}$ for $u\ge0$ and some $a>0$.
The final assumption is the strict potential
\begin{equation}\label{ac-potential}
 \log\E e^{\phi(\sqrt q Z)^2}\le q-cq^2,
                  \qquad 0\le q\le2,
\end{equation}
where $Z$ is standard Gaussian. The main text verifies all these
conditions for the activation constructed there.

Set $H=L-1\ge1$. For a finite input set $X$, a controlled critical
path starts at $Q_0(x,y)=(1+x\cdot y)/d$ and consists of
\[
 \rho_j\ll\mathcal N(0,Q_{j-1}),\quad
 D_j=\operatorname{KL}(\rho_j\Vert\mathcal N(0,Q_{j-1})),\quad
 Q_j(x,y)=\E_{\rho_j}\phi(G_x)\phi(G_y).
\]
Singular references are permitted. Its output $z\in\operatorname{range}Q_H$
has readout cost $r=z^{\mathsf T}Q_H^\dagger z/2$, and total action
$A=\sum_jD_j+r$. Let $J_L(f)$ be the infimum of this action over
full-cube paths with $f(x)z(x)\ge1/2$ everywhere. These definitions
will be related to actual finite-width probabilities below.

\subsection{Scalar control, cooling, and a target program}

For $q>0$ let $V(q)$ be the infimum, over finite scalar paths starting
at $q_0=q$, of
\[
 \sum_{j=1}^m\operatorname{KL}(\eta_j\Vert\mathcal N(0,q_{j-1}))
       +\frac1{8q_m},\qquad q_j=\E_{\eta_j}\phi(G)^2.
\]
The length-zero path is included.

\begin{lemma}[Scalar baseline and short paths]\label{ac-scalar}
The function $V$ is nonincreasing and has a finite positive limit
$V_*:=\lim_{q\downarrow0}V(q)$. For a fixed $C$,
\[
 0\le V_*-V(q)\le Cq\quad(0<q\le2).
\]
A scalar transition of entropy $D$ from variance $q$ to $q'$ satisfies
\[
 D\ge V(q)-V(q'),\qquad D\ge q'-q+cq^2\quad(q\le2).
\]
For every sufficiently small $p>0$ there is a strictly increasing
scalar path starting at $p$, of total cost at most $V(p)+p$, whose
length is at most $C/p$ and for which
\begin{equation}\label{ac-short-path}
 \sum_j q_j\le C\log(2/p),\qquad
 \sum_j\sqrt{D_j}\le C\log(2/p).
\end{equation}
\end{lemma}
\begin{proof}
Rescaling a controlled preactivation when its reference variance
increases preserves its Gaussian relative entropy and increases its
$\phi$-square moment. Applying this successively proves monotonicity.
Oddness, $\phi'(0)=1$, and bounded third derivatives give
$F(u):=\E\phi(\sqrt uZ)^2\ge u-Cu^2$ for small $u$; outside a
fixed neighborhood of zero use Gaussian tails and $|\phi(z)|\le|z|$.
Inflating the reference variance by $1+Kq$, with fixed large $K$,
therefore increases a small $q$ by between $c_1q^2$ and $C_1q^2$
at Gaussian cost at most $C_2q^2$. Continuous adjustment of the
inflation hits any prescribed small endpoint $p$ exactly. Telescoping
the cost-to-increment ratio gives a bridge from arbitrarily small
positive variance to $p$ costing at most $Cp$. A direct readout at
one fixed $p$ proves finiteness of $V_*$, and appending a nearly
optimal suffix proves $V_*-V(p)\le Cp$. Boundedness extends this
inequality to $q\le2$. For starting $q\le1$, every final variance,
including that of a length-zero path, is at most one; hence $V_*\ge1/8$.
The Bellman inequality follows by appending a nearly optimal suffix,
and entropy duality applied to~\eqref{ac-potential} gives dissipation.

Delete each nonincreasing transition of a near-optimal path and rescale
the subsequent controls as above. Cost does not increase, so the path
may be taken strictly increasing. Its cost is bounded by a fixed
constant $B$. Gaussian entropy duality and $\phi(z)^2\le z^2$ bound
each ratio $q'/q$ by a constant depending on $B$. The entropy spent
before a state $q$ is at most $p+Cq$, by the scalar suffix bound.
In a dyadic band $u\le q<2u$, its expenditure, including its exit
transition, is therefore $O(u)$: the exit state is at most a fixed
multiple of $u$. Summed dissipation bounds its length by $O(1/u)$,
apart from at most one exit transition. The band contributes $O(1)$
to both sums in~\eqref{ac-short-path}, using Cauchy--Schwarz for the
second. Summing the geometric length bounds and logarithmically
many bands proves the assertion. States above a fixed small constant
occupy only a bounded number of transitions by the same dissipation.
\end{proof}

\begin{lemma}[An exact low-cost Fourier seed]\label{ac-seed}
Let $A_F(f)=\sum_T|\widehat f(T)|$. There is a first-row controlled
law, with auxiliaries and a score $b_1$, such that
\[
 \E b_1^2=1,\quad \E[b_1\phi(w\cdot(1,x))]=\sqrt{v_1}f(x),
 \quad \E\phi(w\cdot(1,x))^2=q_1
\]
with $q_1$ independent of $x$, relative entropy at most $1/H$, and
$v_1^{-1},\norm{b_1}_\infty$ bounded by a fixed polynomial in
$(d+H+2)A_F(f)$. The reference weight marginal is exactly
$\mathcal N(0,I/d)$.
\end{lemma}
\begin{proof}
For centered Gaussian $G$ of variance at most two, split at $G=-M/2$
to obtain
\[
 \beta-\E\phi(M+G)\le CM^{-a}+Ce^{-M^2/16}.
\]
A parity on $k$ coordinates is a signed sum of $k$ threshold steps
in their Hamming sum and a constant, with coefficient magnitudes
at most one. Each step has an affine argument separated from zero
by one. Shift first-row Gaussian means by a polynomially large
multiple of these affine coefficients, and mix the signed detectors;
reflection of the row implements a negative sign. Their normalized
mixture has mean $g$ with $\chi_T(x)g(x)\ge c/(k+1)$ everywhere.
Every component costs at most $D_* =\poly(d)$ in Gaussian KL.

Average this law over cube translations $x\mapsto s x$, reflecting
the output row by $\chi_T(s)$. Its mean becomes exactly
$c_T\chi_T(x)$, where $c_T=\E_x\chi_T(x)g(x)\ge c/(d+1)$.
Its second-moment diagonal is constant. The transformations preserve
the Gaussian reference, and entropy convexity preserves the cost
$D_*$. Mix these laws, with reflections according to
$\operatorname{sign}\widehat f(T)$, in proportions
$|\widehat f(T)|/c_T$. The resulting law has mean $f/Z_f$,
constant diagonal, cost at most $D_*$, and
$Z_f\le C(d+1)A_F(f)$.

Dilute this law into the free first-row law with probability
$\theta=(HD_*)^{-1}$, increasing $D_*\ge1$ if needed. Retain its
Bernoulli tag $B_0$ and use $b_1=B_0/\sqrt\theta$. The coefficient
is $\sqrt\theta f/Z_f$, so $v_1=\theta/Z_f^2$. In the reference
use the same tags and free Gaussian weights in every branch. Its
weight marginal is the original Gaussian, and the joint entropy is
at most $\theta D_*\le1/H$. All claimed bounds follow. In particular,
no entropy proportional to the number of translation or Fourier tags
has been charged.
\end{proof}

\begin{lemma}[A scalar-channel target program]\label{ac-program}
Put $R=(d+L+2)A_F(f)$ and $S=1+\log R$. If $H\ge CS$, there is
an exact $H$-layer controlled program with scores $b_j$ and deterministic
positive components $v_j$, residual diagonals $r_j\ge0$, such that
\[
 \E b_j^2=1,\quad \E[b_jh_j(x)]=\sqrt{v_j}f(x),\quad
 \E h_j(x)^2=v_j+r_j\le1.
\]
Its charged scalar costs $D_j$ obey
\begin{equation}\label{ac-program-bounds}
 A:=\sum_jD_j+\frac1{8v_H}\le V_*+C S^2/H,
 \quad K:=\sum_{j=2}^H\sqrt{D_j}\le CS,
 \quad |b_j|\le B\le e^{CS}.
\end{equation}
For $j\ge2$, the controlled row uses a scalar $Z$, of entropy cost
at most $D_j$ relative to a reference with standard Gaussian $Z$
marginal, and independent Gaussian residual noise. At one input its
feature is
$\phi(\sqrt{v_{j-1}}f(x)Z+\sqrt{r_{j-1}}G)$.
The score depends on $Z$ and auxiliaries, not on $G$.
\end{lemma}
\begin{proof}
We supply the scalar estimates and the depth accounting. Write
$b(u)=\E\phi'(\sqrt uZ)$. Taylor expansion, with bounded fifth
derivatives and Gaussian tails, gives
\begin{equation}\label{ac-cooling-expansions}
 F(u)=u+\phi'''(0)u^2+O(u^3),\quad
 b(u)=1+\tfrac12\phi'''(0)u+O(u^2),\quad
 \frac{u b(u)^2}{F(u)}=1+O(u^2).
\end{equation}
Also $F(u)\le u-cu^2$ by~\eqref{ac-potential}. In a free layer,
Gaussian integration by parts gives at least $v b(q)^2$ as the
component generated from $vff^{\mathsf T}$ at diagonal $q$.
The component-to-diagonal ratio therefore loses factors
$q b(q)^2/F(q)$. Their logarithms are $O(q^2)$ near zero, and
they have a positive lower bound on each remaining compact interval.
Since free dissipation bounds $\sum q_j^2$, cooling to the first
diagonal at most $\delta$ takes $O(1/\delta)$ layers and retains
component at least $c\delta v_1$. If already below $\delta$, skip
cooling; the same lower bound holds. Cooling costs zero entropy.

The estimate
\begin{equation}\label{ac-derivative-comparison}
 1-\phi'(z)\le C\phi(z)^2
\end{equation}
follows near zero by odd $C^3$ regularity and away from zero by
positivity of $\phi$. Inflate only the target Gaussian coordinate
by variance factor $1+e$. Its cost is
$[e-\log(1+e)]/2\le e^2/4$. If the current component and residual
diagonal are $v,r$ and $q=v+r$, the full conditional-mean component
and new residual satisfy
\begin{equation}\label{ac-growth}
 v'\ge(1+e)v b(q+ev)^2
       \ge(1+e)v[1-C(q+ev)],\qquad r'\le r
\end{equation}
in the small-variance range. The first inequality is the first Hermite
projection; the last follows from conditional Gaussian Poincare.
The component is continuous and increasing in $e$, and at $e=0$
it is at most $v$. Thus at $q\le\delta$ one can hold $v$ exactly
with $e=O(\delta)$, at cost $O(\delta^2)$, preserving $q\le\delta$.
If $r\le\delta$ and $v\le p$, one can prescribe
\[
 v' =\min\{p,v[1+c_1(\delta+v)]\}
\]
with $e=O(\delta+v)$ and cost $O((\delta+v)^2)$. From
$v_0=\delta e^{-C_1S}$ this takes $O(S/\delta)$ steps and costs
$O(p+S\delta)$. Indeed, on full steps,
$(\delta+v)^2=O(\Delta v+\delta\Delta v/v)$;
their square-root costs sum to $O(\log(p/v_0)+1)=O(S)$ when
$p/\delta$ is fixed. The prescribed sequence fixes the length
independently of the residual and any preceding holding time.

At $v=p$, $r\le\delta\le p$, hold the full component exactly.
The total Gaussian preactivation variance lies in $[p,Cp]$, and
\[
 \E\phi'(\sqrt uZ)^2=1+\phi'''(0)u+O(u^2)\le1-cp.
\]
The derivative form of conditional Poincare now gives
$r'\le(1-cp)r$. Thus $O(\log(H+2)/p)$ purification stages leave
$r\le\delta H^{-4}$, at cost $O(p\log(H+2))$ and square-root
cost $O(\log(H+2))$.

For the scalar suffix, $|\phi''(z)|\le C\phi(z)$ for $z>0$:
use oddness and $C^3$ regularity near zero, positivity on compact
intervals, and bounded $\phi''$ elsewhere. If $P_r$ denotes Gaussian
heat convolution, then
\begin{equation}\label{ac-heat}
 e^{-Cr/2}\phi(z)\le P_r\phi(z)\le\phi(z),\qquad z\ge0.
\end{equation}
Indeed $g=\phi''+C\phi$ is odd and nonnegative on the positive
half-line. Pairing the Gaussian density at $y,-y$ shows
$P_sg(z)\ge0$ for $z\ge0$, so
$\partial_s[e^{Cs/2}P_s\phi(z)]\ge0$. The upper bound follows
in the same way from concavity and the odd function $\phi''$.
Concavity also gives $\phi(\alpha z)\ge\alpha\phi(z)$ for
$0\le\alpha\le1$ and $z\ge0$.

Lift the short scalar path of Lemma~\ref{ac-scalar} starting at $p$.
For a nominal law $\eta$ from variance $q$ and an actual component
$v\le q$, use physical scalar $\sqrt{v/q}X$, $X\sim\eta$.
Its Gaussian relative entropy is unchanged. Equation~\eqref{ac-heat}
and concavity imply
\[
 e^{-Cr}(v/q)\le v'/q'\le1,
\]
and the residual remains at most $r$. Starting at $v=q=p$, the
lift of $m\le C/p$ stages therefore retains at least
$e^{-Cm\delta H^{-4}}$ of every nominal component. Since
$\delta\le p$, all components stay at least $p/2$, and the extra
readout cost is $O(H^{-4})$. Its charged square-root cost is
$O(\log(2/p))$ by~\eqref{ac-short-path}.

Choose $\delta=B_1S/H$ and $p=B_2S/H$, with fixed sufficiently
large $B_2\ge B_1$, and then increase the constant in $H\ge CS$
so both scales are small. After cooling reduce the component to
$v_0=\delta e^{-C_1S}$, which is permitted by Lemma~\ref{ac-seed}.
Cooling, prescribed growth, purification, and the chosen scalar
suffix occupy at most $H-1$ stages after choosing $B_1,B_2$.
Insert exactly the unused number of holding stages before growth.
The total cost is
\[
 A\le H^{-1}+V(p)+p+
 C[H\delta^2+p+S\delta+p\log(H+2)+H^{-4}]
 \le V_*+CS^2/H.
\]
Holding contributes $O(H\delta)=O(S)$ to the square-root sum;
the other contributions were bounded above, proving $K\le CS$.
This is a bound on the charged costs themselves, not an inference
from smaller marginal entropies.

The full conditional-mean score is
$P_r\phi(\sqrt vZ)/\sqrt{v'}$, of unit second moment. To retain
a smaller prescribed component $0<v''\le v'$, replace it by
\[
 \sqrt{v''/v'}\,b+\sqrt{1-v''/v'}\,U,
\]
where $U$ is an independent random sign with the same reference law.
This has unit second moment, the desired coefficient, and absolute
bound $1+(v'')^{-1/2}$, with no extra entropy. Every retained component
in the construction is at least $e^{-CS}$, proving the score bound.
At each stage an independent Gaussian residual vector with the current
residual covariance realizes the stated full-cube law. Only its
constant diagonal enters the one-input identities.
\end{proof}

\subsection{Finite-width realization without a rate-limit argument}

\begin{lemma}[Finite scalar-channel probability]\label{ac-channel}
Suppose a program as in Lemma~\ref{ac-program} has action $A$,
distortion $K=\sum_{j\ge2}\sqrt{D_j}$, and score bound $B\ge1$.
For every $\tau>0$ and every integer
\begin{equation}\label{ac-channel-width}
 N\ge CB^4H^4e^{4K}(A+1)^6(1+\tau^{-1})^4
       [\log(2|X|)+\log(H+A+\tau^{-1}+3)],
\end{equation}
the original critical Gaussian prior satisfies
\[
 \mu_N\{f(x)z_W(x)>1/2\ (x\in X)\}\ge e^{-N(A+\tau)}.
\]
The lemma uses only boundedness and the unit Lipschitz constant of the
activation, in addition to the specified program.
\end{lemma}
\begin{proof}
For $N$ independent score/feature rows set
\[
 M=N^{-1}\sum_i b_i^2,\quad
 \widehat c(x)=\frac{N^{-1}\sum_i b_ih_i(x)}{\sqrt M},\quad
 \widehat q(x)=N^{-1}\sum_i h_i(x)^2,\quad
 \widehat s(x)=\sqrt{\widehat q(x)-\widehat c(x)^2}.
\]
If the three raw moments have errors at most $\alpha\le1/4$,
then, writing $c=\E bh$ and $s=\sqrt{\E h^2-c^2}$,
\begin{equation}\label{ac-score-metric}
 \max_x\sqrt{(\widehat c-c)^2+(\widehat s-s)^2}
                         \le8\sqrt\alpha.
\end{equation}
For example $|\widehat c-c|\le4\alpha$ and the squared residual
error before taking square roots is at most $13\alpha$.
The coefficients $a_i=b_i/(N\sqrt M)$ have norm squared $1/N$;
the empirical kernel decomposes exactly as
$K_{\rm emp}=\widehat c\widehat c^{\mathsf T}+R_{\rm emp}$ with
$R_{\rm emp}\succeq0$.

Condition on the entire preceding construction and set $u=\sqrt N a$.
For a fresh row draw its prescribed scalar $(Z,U)$ and independent
standard Gaussian $\xi\in\R^N$, and use
\begin{equation}\label{ac-row-realization}
 w=[Zu+(I-uu^{\mathsf T})\xi]/\sqrt N.
\end{equation}
Under the reference, $Z$ is standard Gaussian, so the incoming
weight marginal is exactly $\mathcal N(0,I/N)$. Its controlled
relative entropy is at most $D_j$ by data processing. At one input
the preactivation has law $\widehat c(x)Z+\widehat s(x)G$,
with independent standard $G$. Couple this to the ideal preactivation
using the same $Z,U,G$. Gaussian entropy duality gives
$\E Z^2\le e^{2\sqrt{D_j}}$, so the preactivation and feature RMS
errors are at most $e^{\sqrt{D_j}}$ times the preceding metric error.
Orthogonal projection onto the common unit score, followed by the
reverse triangle inequality for residual norms, gives the same bound
for the population coefficient/residual metric. Fresh sampling and
\eqref{ac-score-metric} consequently give
\begin{equation}\label{ac-metric-recursion}
 e_H\le8He^K\sqrt\alpha.
\end{equation}
Only one-input comparison couplings are used; the sampled network
keeps its usual joint residual Gaussian law on all inputs.

The observables $b^2,bh(x),h(x)^2$ are bounded by $B^2,B,1$.
Conditional Hoeffding, a union bound over $2|X|+1$ observables, and
successive conditioning over layers give all moment estimates with
probability at least $1-\kappa/2$ if
\[
 N\ge CB^4\alpha^{-2}\log[8H(2|X|+1)/\kappa].
\]
Stop controlling if $M\notin[1/2,3/2]$ and thereafter use the prior.
This defines the law on all histories, while the good event never stops.
Choose
\[
 \kappa=\min\{1/4,\tau/[32(A+1)]\},\qquad
 \alpha=\left[\frac{\kappa\sqrt{v_H}}{64He^K}\right]^2.
\]
Then the final coefficient has signed margin at least
$\sqrt{v_H}(1-\kappa/8)$. On this event shift the readout mean to
$ta$, where $t=(1+\kappa)/(2\sqrt{v_H})$; otherwise use zero mean.
Its mean output is at least $1/2+\kappa/4$ in the correct direction.
Its noise variance is at most $1/N$ at each input. If
$N\ge C\kappa^{-2}\log(4|X|/\kappa)$, the maximum noise is less
than $\kappa/8$ with probability at least $1-\kappa/2$. Thus the
controlled law $Q_N$ gives the strict target event probability at
least $1-\kappa$.

Use the same reference auxiliaries as in the program. The reference
weight marginal in~\eqref{ac-row-realization} remains Gaussian even
under the history-dependent orthogonal decomposition. Conditional
chain rule and data processing give hidden-block cost at most $ND_j$;
stopping only reduces this bound. The readout shift costs $Nt^2/2$.
After discarding auxiliaries,
\[
 \operatorname{KL}(Q_N\Vert\mu_N)
 \le N\left[\sum_jD_j+\frac{(1+\kappa)^2}{8v_H}\right]
 \le N(A+3A\kappa).
\]
Binary relative entropy for the target event implies
\[
 -N^{-1}\log\mu_N(E_f)
 \le\frac{A+3A\kappa+N^{-1}\log2}{1-\kappa}\le A+\tau
\]
when $N\ge8/\tau$. Since $v_H^{-1}\le8A$, the width bound
\eqref{ac-channel-width} implies all the numerical conditions above.
\end{proof}

\begin{theorem}[Critical raw target mass]\label{ac-target-mass}
For fixed $\phi$ satisfying the hypotheses of this appendix, there
are constants $C_0,C_1,C_2$ such that, if
$R=(d+L+2)A_F(f)$ and $L\ge C_0\log R$, then at every integer
$N\ge R^{C_1}$,
\[
 \mu_N\{f(x)z_W(x)>1/2\ (x\in X_d)\}
 \ge\exp\left[-N\left(V_*+C_2\frac{\log^2R}{L}\right)\right].
\]
\end{theorem}
\begin{proof}
Apply Lemma~\ref{ac-channel} to Lemma~\ref{ac-program} with
$\tau=1/H$. Each factor in~\eqref{ac-channel-width} is at most
$e^{CS}$: $B\le e^{CS}$, $K\le CS$, $H,d\le e^S$, and
$A+1\le C(1+S)$. Also $\log|X_d|=d\log2$. The additional
$1/H$ action is absorbed into $CS^2/H$. Enlarging constants handles
$H=L-1$ and $S=1+\log R$.
\end{proof}

\subsection{Near-baseline paths have parallel circuits}

For $D\ge0$ define $G(D)\ge1$ by
$G(D)-1-\log G(D)=2D$. Gaussian entropy duality implies
\begin{equation}\label{ac-gaussian-gain}
 \E_\rho\langle g,u\rangle^2\le G(D)\norm u^2,
 \qquad \sqrt{G(D)}\le e^{\sqrt D},\qquad G(D)\le4D+2
\end{equation}
for any tilt of a standard Gaussian with entropy at most $D$.
Indeed projection onto $\langle g,u\rangle/\norm u$ gives
$D\ge(m-1-\log m)/2$ when its second moment $m>1$; writing
$m=e^{2v}$ gives $(m-1-\log m)/2\ge v^2$.

\begin{lemma}[Variance and entropy budgets]\label{ac-occupation}
For a full-cube globally $1/2$-margined path of action
$A\le V_*+\epsilon$, $\epsilon\ge0$, put
$q_j=\max_xQ_j(x,x)$. Then
\[
 \sum_{j=0}^{H-1}q_j+\sum_{j=1}^H\sqrt{D_j}
          \le C[1+\log(H+2)+\sqrt{H\epsilon}].
\]
\end{lemma}
\begin{proof}
All one-input variances are positive, since a zero variance would
propagate at finite entropy and force zero final output. Decrease $c$
so $q-cq^2$ is increasing on $[0,2]$. Projecting onto an input attaining
the new maximum and using Lemma~\ref{ac-scalar} yields
\[
 D_j\ge q_j-q_{j-1}+cq_{j-1}^2.
\]
The scalar suffix at an input attaining $q_i$, including its readout,
costs at least $V(q_i)$. Thus
\begin{equation}\label{ac-prefix-budget}
 \sum_{j\le i}D_j\le\epsilon+Cq_i.
\end{equation}
These inequalities concern full layer entropies. No two-input
contraction property of the activation is needed.

For an interval of $t$ transitions whose pre-states and endpoint
are at least $u$, dividing dissipation by the squared pre-state and
using $(q'-q)/q^2\ge1/q-1/q'$ gives
$t\le C(E/u^2+1/u)$, where $E$ is its entropy expenditure.
An initial prefix of length $t$ and expenditure $E$ also obeys
\begin{equation}\label{ac-forced-cooling}
 \sum_{j<t}q_j\le C\log(t+2)+C\sqrt{tE}
\end{equation}
when total entropy is bounded by a fixed constant. To prove this,
write $F_0(v)=v-cv^2$, compare $q_j$ to $r_j+s_j$ with
$r_{j+1}=F_0(r_j)$, $s_{j+1}=F_0(s_j)+D_{j+1}$, $r_0=q_0\le2$,
$s_0=0$, and decrease $c$ so $F_0$ is increasing on the resulting
bounded range. Since $F_0(r+s)\le F_0(r)+F_0(s)$, induction proves
the comparison. The free sum is $O(\log(t+2))$, and
$c\sum s_j^2\le E$, proving~\eqref{ac-forced-cooling}.

If $\epsilon$ exceeds a fixed positive constant, summed dissipation
and Cauchy--Schwarz prove the claim. Otherwise total entropy is bounded.
Set $a=\max\{\epsilon,(H+1)^{-1}\}$, $b=\min_jq_j$, and
$u_0=\max\{a,b\}$. If $u_0$ exceeds a fixed small constant, the
occupation bound gives bounded $H$. Otherwise use dyadic levels
$u_j=2^ju_0$ until a fixed constant, and their last-exit indices
$\tau_j=\max\{i:q_i\le u_j\}$. There are $O(\log(H+2))$ bands.
The initial prefix has entropy $O(\epsilon+1/H)$ if $b\le a$;
then~\eqref{ac-forced-cooling} and Cauchy--Schwarz bound its two
contributions as claimed. If $b>a$, its entropy is $O(b)$ and its
length $O(1/b)$, giving state sum $O(\log(H+2))$ and square-root
entropy $O(1)$. A later nonempty last-exit block has entropy $O(u_j)$
by~\eqref{ac-prefix-budget}; after removing its first transition
its pre-states and endpoint exceed $u_j$, so its length is $O(1/u_j)$.
Every state in it is $O(u_j)$ because its starting state is at most
$u_j$ and any increase is at most spent entropy. Both contributions
are $O(1)$. The final tail has bounded length and contributions.
Summing proves the assertion.
\end{proof}

\begin{theorem}[Critical canonical circuit theorem]\label{ac-circuits}
A globally $1/2$-margined path of action $A\le V_*+\epsilon$ gives
a nonuniform threshold circuit of size $R_1^{O(1)}$ and depth
$O(\log R_1)$, and a bounded-fan-in Boolean circuit of size
$R_1^{O(1)}$ and depth $O(\log^2R_1)$, where
\[
 T=1+\log(H+2)+\sqrt{H\epsilon},\qquad
 R_1=(d+H+A+2)e^{CT}.
\]
All constants depend only on the fixed activation.
\end{theorem}
\begin{proof}
We give the anchor reduction, finite-dimensional state, evaluator,
and parallel iteration explicitly.

\paragraph{Anchor reduction.}
Let $I_j(U)$ be the entropy of the layer-$j$ preactivation marginal
on an anchor set $U$. It increases with $U$ and is at most $D_j$.
Choose $\rho=[C H(A+1)]^{-1}\le1/64$ and a cutoff $\delta>0$.
Greedily add a query whenever $I_j(U\cup\{x\})-I_j(U)>\delta$
for some layer, or projection of the final readout onto the anchor
feature span changes its output there by more than $\rho$.
There are at most $A/\delta$ entropy additions and $2r/\rho^2$
readout additions: each latter increases the projection's squared
norm by at least $\rho^2$, since final feature norms are at most one.

Replace each row density by its conditional expectation $w_j(G_U)$
given the anchors. It has entropy $I_j(U)\le D_j$ and preserves
the original anchor marginal exactly. The conditional query KL
increment is at most $\delta$. Conditional maximal coupling and
Pinsker give expected squared feature error at most
$4\sqrt{\delta/2}$, since features are bounded by one.
For kernels sharing their anchor block, compare a query by its
projection onto the common anchor span and its orthogonal residual
norm. This defines a metric $d_U$ equal to the RMS error in a Gaussian
coupling sharing the anchors and one residual standard Gaussian.
The preceding replacement has $d_U$ error $O(\delta^{1/4})$.
For a later Gaussian reference with the same anchor block, tilt this
one-query coupling by the common $w_j$. Its entropy is $I_j(U)$;
\eqref{ac-gaussian-gain} and the Lipschitz property multiply the
input error by at most $e^{\sqrt{D_j}}$. Thus the complete error is
at most $CH\delta^{1/4}\exp(\sum_j\sqrt{D_j})$.
Choose $\delta$ so this is at most
$\rho/\sqrt{\max\{2r,1\}}$. Lemma~\ref{ac-occupation} makes
the anchor count polynomial in $R_1$.

These replacements define a \emph{genuine} full-cube controlled path,
with the exact original anchor blocks, density depending only on
anchors, and no greater row entropies. The projected readout loses
at most $2\rho$ of margin. Scaling it by $(1-4\rho)^{-1}$ restores
margin $1/2$ at action increment $O(1/H)$. Its coefficient vector
has norm $O(\sqrt{A+1})$. The genuine path therefore satisfies
the budgets of Lemma~\ref{ac-occupation} with bound $CT$.

\paragraph{Canonical finite states.}
Initially $u_0(x)=(1,x)/\sqrt d\in\R^{d+1}$. Given states in
$\R^K$, let $g\sim\mathcal N(0,I_K)$ tilted by
$w_j((\langle g,u(a)\rangle)_{a\in U})$, and write this law as
$\rho_j$. In $L^2(\rho_j)$ put
$\Phi(u)=\phi(\langle g,u\rangle)$ and let $V$ span all $K$
coordinate functions $g_i$ and all anchor features $\Phi(u(a))$.
For an orthonormal basis $e_1,\ldots,e_M$, $M\le K+|U|$, define
\begin{equation}\label{ac-canonical}
 F(u)=\bigl(\langle e_b,\Phi(u)\rangle_{b=1}^M,
                         \norm{(I-P_V)\Phi(u)}\bigr).
\end{equation}
Pythagoras preserves every diagonal. At anchors the last coordinate
is zero, so every query-anchor cross moment and the entire anchor
block are preserved exactly. Future anchor-plus-one-query Gaussian
laws are determined by precisely these moments. Induction gives
$d_U=0$ relative to the genuine anchor-only path at every stage,
and identical anchored readout. The canonical full covariance is
not asserted to be a valid original transition. Its state dimension
is at most $d+1+H(|U|+1)$, polynomial in $R_1$.

Since the linear map $u\mapsto\langle g,u\rangle$ lies in $V$,
write $F(u)=A_ju+B_j(u)$, with
$\norm{A_j}_{\rm op}\le\alpha_j:=\sqrt{G(D_j)}$.
The map taking a Hilbert residual to its $V$ projection and its
orthogonal norm is $1$-Lipschitz. Moreover
$|\phi'(s)-1|\le C\min(s^2,1)$ by odd $C^3$ regularity and
$0\le\phi'\le1$. For $\norm u,\norm v\le\sqrt q$, put
$X=g\cdot u$, $Y=g\cdot v$, $Z=g\cdot(u-v)/\norm{u-v}$.
The squared normalized residual difference is at most a fixed
multiple of
$W=Z^2\min\{1,\max(|X|,|Y|)^4\}$.
For small fixed $c>0$,
\[
 e^{cW}-1\le cZ^2(X^4+Y^4)e^{cZ^2},\qquad
 \log\E_{\rm Gauss}e^{cW}\le Cq^2.
\]
For the second inequality condition $X$ or $Y$ on $Z$; its fourth
moment is at most $Cq^2(1+Z^4)$, regardless of correlations.
Entropy duality consequently gives
\begin{equation}\label{ac-residual-budget}
 \Lip(B_j)\le C(q+\sqrt{D_j}),\quad
 \sum_j\Lip(B_j)\le CT,\quad \prod_j\alpha_j\le e^{CT}
\end{equation}
on the actual state balls. Here use the genuine path's diagonal
budget and the original path's entropy budget; alternatively take
$D_j$ to be the actual anchor entropy. Extend each $B_j$ by radial
projection. Round the radius upwards on mesh $R_1^{-c_1}$ and impose
that same lower floor, with fixed large $c_1$. This adds only a
constant to the summed squared radii and preserves the actual path.

\paragraph{Polynomial finite evaluation.}
For any current law $\rho$, basis vector $e=(e_b)$, and polynomial
input radius, write
\[
 a(u)=\E_\rho e(g)\phi(g\cdot u),\quad
 b(u)=\E_\rho\phi(g\cdot u)^2,\quad
 F(u)=(a(u),\sqrt{b(u)-\norm{a(u)}^2}).
\]
Equation~\eqref{ac-gaussian-gain} gives
$\E_\rho\norm g^2\le K G(D)$. Truncate to $\norm g\le M_0$.
If the discarded probability is $p$, Bessel's inequality bounds the
discarded vector expectation uniformly by $\sqrt p$, and the scalar
expectation by $p$. A polynomial $M_0$ gives any inverse-polynomial
accuracy. On the retained set sample with importance density
$(1+\norm e^2)/Z_0$ relative to $\rho$, normalized there, where
$Z_0\le M+1$. The vector and scalar integrands have norms at most
$Z_0/2,Z_0$, and Lipschitz bounds at most $Z_0M_0/2,2Z_0M_0$.
An input net has logarithmic cardinality
$O(K\log(2+\poly(R_1)/\theta))$. For bounded Hilbert-valued
samples, expected mean error is at most the bound divided by the
square root of the sample count; bounded differences adds an
exponential tail. A union bound on this net therefore gives
polynomially many nodes approximating $a,b$ uniformly to error
$\theta$. Lipschitzness extends the estimate to the whole ball.
The retained nodes and all importance-weighted coefficients have
polynomial magnitude, irrespective of an anchor condition number.
Since $\norm a\le1$, choosing $\theta=O(\eta^2)$ and clipping
the approximate radicand at zero gives error $O(\eta)$ in $F$.

Round these constants, the matrix $A_j$, and all arithmetic to
inverse-polynomial accuracy. Scalar $\phi$ values need only a
polynomial mesh, since $\phi$ is $1$-Lipschitz; no algorithm for
evaluating $\phi$ is assumed. Every word has $O(\log R_1)$ bits.
Unary and binary word operations have polynomial-size truth tables.
A sum of polynomially many such words has polynomial integer range
after a common fixed-point scaling; threshold tests against each
possible integer and a bit decoder compute it in constant threshold
depth. Signed integer weights of polynomial magnitude can be
replaced by repeated inputs and complements. Squares, square roots,
comparisons, clipping, and radial projection have the same bounds;
division occurs only above the inverse-polynomial radius floor.
Thus one canonical residual evaluator has polynomial size and
constant threshold depth, uniformly on the required polynomial ball.

\paragraph{Parallel iteration.}
Set $S_j=\prod_{i\le j}\alpha_i$, $v_j=u_j/S_j$. The normalized
recursion is $v_j=C_jv_{j-1}+H_j(v_{j-1})$, where
$C_j=A_j/\alpha_j$ is a contraction and
$\Lip(H_j)\le\ell_j:=C(q_{j-1}+\sqrt{D_j})$ with
$\Lambda:=\sum_j\ell_j\le CT$. All $S_j$ are polynomial in $R_1$.
Precompute the interval products $C_{j:i}=C_j\cdots C_i$ as
nonuniform constant matrices; their operator norms and entries are
at most one. Starting with $v_j^{(0)}=C_{j:1}v_0$, compute all times
simultaneously by
\[
 v_j^{(k+1)}=C_{j:1}v_0+
               \sum_{i=1}^j C_{j:i+1}H_i(v_{i-1}^{(k)}).
\]
Strict time ordering gives
$\max_j\norm{v_j^{(k)}-v_j}\le2\sqrt2\Lambda^k/k!$:
expansion uses distinct ordered indices, whose products sum to at
most $\Lambda^k/k!$. Thus $O(\log R_1)$ iterations give any fixed
inverse-polynomial accuracy. Uniform additive error $\delta_0$ per
update is amplified by at most
$\sum_{k\ge0}\Lambda^k/k!=e^\Lambda=R_1^{O(1)}$.
The clipped residuals have bounded ranges, so all exact and inexact
iterates stay in a fixed polynomial ball. Rounded interval products,
quadrature, and arithmetic contribute such additive errors; divide
primitive tolerances by a polynomial to absorb the number of terms.

There are polynomially many coordinates and times, each update has
constant threshold depth, and there are $O(\log R_1)$ updates.
The final rescaling and readout amplify errors only polynomially;
choose precision so the output error is below $1/16$. This preserves
all signs. Replacing each polynomial-fan-in majority gate by a
logarithmic-depth Boolean summation circuit proves the Boolean bound.
\end{proof}

\subsection{A finite-event certificate and the full-cube ceiling}

\begin{lemma}[Bounded-dual extraction]\label{ac-dual}
On $m$ inputs put $k=m(m+1)/2$. For any raw event $E$ of critical
prior probability at least $e^{-NA}$, any $\tau>0$, and
$0<\Delta\le1$, there is a realization in $E$ and laws relative
to its empirical preceding Gaussian covariances such that
\[
 \max_{x,y}|\E_{\rho_j}\phi(G_x)\phi(G_y)-Q_j(x,y)|\le\Delta,
 \qquad \sum_jD_j+r\le A+\tau,
\]
provided
\begin{equation}\label{ac-dual-width}
 N\ge\frac C\tau\left[
 Hk\log\left(1+\frac{C(A/\tau+1)Hk}{\Delta}\right)
 +m\log(C(A/\tau+1))+1\right].
\end{equation}
Here $z$ is the realization's actual output and
$r=z^{\mathsf T}Q_H^\dagger z/2$. No symmetry of the event is needed.
\end{lemma}
\begin{proof}
For a reference feature law $\mu$ on $[-1,1]^m$, let $\Phi(h)$ list
its $k$ upper-triangular products and define
\[
 I_B(q)=\max_{\lambda\in[-B,B]^k}
       [\lambda\cdot q-\log\E_\mu e^{\lambda\cdot\Phi}].
\]
A maximizer exists. Its exponential tilt $\nu$ has moment vector
$a$; coordinatewise optimality on the box gives
$\lambda_i(q_i-a_i)=B|q_i-a_i|$. Hence
\begin{equation}\label{ac-dual-identity}
 I_B(q)=\operatorname{KL}(\nu\Vert\mu)+B\norm{q-a}_1\ge0.
\end{equation}
The same tilt of Gaussian preactivations has identical entropy and
moments. The objective is $2k$-Lipschitz in the sup norm of $\lambda$.
A mesh-$\eta$ grid has at most $(1+CB/\eta)^k$ points, and each
fixed exponential likelihood ratio for $N$ iid rows has mean one.
Thus for empirical moments $\widehat q$,
\[
 \E e^{N I_B(\widehat q)}\le e^{2Nk\eta}(1+CB/\eta)^k.
\]
This is uniform in the conditional reference. Iterate through layers.
Given the hidden layers, the actual readout covariance is $Q_H/N$,
so $\E[e^{(1-s)Nr}\mid\text{hidden}]\le s^{-m/2}$ for $0<s<1$.
Writing $I_j$ for the realized dual costs gives
\[
 \E e^{N[\sum I_j+(1-s)r]}
 \le e^{2NHk\eta}(1+CB/\eta)^{Hk}s^{-m/2}.
\]
Conditioning on $E$ has relative entropy at most $NA$. Entropy
duality bounds the conditional mean of $\sum I_j+(1-s)r$ by
\[
 A+2Hk\eta+N^{-1}[Hk\log(1+CB/\eta)+(m/2)\log(1/s)].
\]
Take $s=\tau/[16(A+\tau)]$, $\eta=\tau/(64Hk)$, and
$B=(A+\tau)/\Delta$. Under~\eqref{ac-dual-width}, some realization
in $E$ has weighted cost at most $A+\tau/8$, hence unweighted cost
at most $A+\tau$, since all $I_j\ge0$ and
$(A+\tau)(1-s)=A+15\tau/16$. Apply~\eqref{ac-dual-identity}
at each layer; its moment error is at most $I_j/B\le\Delta$.
\end{proof}

\begin{lemma}[Consistent-path repair]\label{ac-repair}
A certificate of total cost at most $B$ in Lemma~\ref{ac-dual}
has a genuine controlled path of cost at most $B$ and output
$\widehat z$ satisfying
\[
 \max_x|\widehat z(x)-z(x)|
            \le\sqrt{2B}\,H\sqrt{m\Delta}\,e^{\sqrt{HB}}.
\]
\end{lemma}
\begin{proof}
Two positive semidefinite $m$-point kernels differing entrywise by
at most $\Delta$ have square-root feature embeddings with row
distances at most $\sqrt{m\Delta}$. Their operator difference is
at most $m\Delta$, and
$\norm{K^{1/2}-{K'}^{1/2}}_{\rm op}
\le\sqrt{\norm{K-K'}_{\rm op}}$.
For the last inequality use $K\le K'+aI$, monotonicity of square
root, and $\sqrt{K'+aI}\le\sqrt{K'}+\sqrt a I$, then reverse
$K,K'$. Existing embeddings can be identified isometrically with
these square-root embeddings, adjoining orthogonal complements.

At each step couple Gaussian rows for the genuine preceding kernel
and the empirical reference in such a common feature space. Tilt
the joint Gaussian by the certificate density on its original
coordinate. The new marginal has no larger entropy. Equation
\eqref{ac-gaussian-gain} and the Lipschitz property multiply each
row RMS error by at most $e^{\sqrt{D_j}}$. Identify the controlled
feature moments with the next empirical kernel using the certificate
error, at additional row distance $\sqrt{m\Delta}$. Gluing the
shared feature spans gives
$e_j\le e^{\sqrt{D_j}}e_{j-1}+\sqrt{m\Delta}$, $e_0=0$.
Consequently $e_H\le H\sqrt{m\Delta}e^{\sqrt{HB}}$.
Transfer the empirical readout vector of norm $\sqrt{2r}$ in the
common space and project it onto the genuine final feature span.
Its norm does not increase, and its output changes by at most
$\sqrt{2r}e_H\le\sqrt{2B}e_H$.
\end{proof}

\begin{theorem}[Actual critical posterior circuit ceiling]\label{ac-ceiling}
Put $R_0=d+n+L+2$, and suppose $L\ge C\log R_0$. At the actual
critical scale $\sigma^2=1/N$, for every integer
$N\ge2^{Cd}R_0^C$ and every consistently labeled sample of at most
$n$ inputs, the original posterior gives probability at least
$31/32$ to a fixed full-cube class of functions having Boolean
circuits of size $\poly(R_0)$. More precisely this class can be taken
as $\{g:J_L(g)\le V_*+C\log^2R_0/L\}$; its threshold depth is
$O(\log R_0)$ and its Boolean depth is $O(\log^2R_0)$.
\end{theorem}
\begin{proof}
The memorizer $g_S=1-2\sum_{p\text{ distinct negative sample input}}
\boldsymbol1_{\{p\}}$ is Boolean, has the sample labels, and has
Fourier variation at most $1+2n$, since each point indicator has
Fourier variation one. Theorem~\ref{ac-target-mass} gives a raw
conditioning denominator at least $e^{-Na_0}$, where
$a_0=V_*+C_0\log^2R_0/L$, at all $N\ge R_0^{C_0}$ after
enlarging constants.

Let $E$ be the raw event of globally margined draws whose Boolean
function satisfies $J_L(g)>a_0+K_0/L$. This is measurable because
there are finitely many Boolean functions on the cube. If
$\mu_N(E)\ge e^{-N(a_0+1/L)}$, apply Lemma~\ref{ac-dual} on
the entire cube with $m=2^d$, action parameter $a_0+1/L$, and
$\tau=1/L$. Choose $\Delta$ so Lemma~\ref{ac-repair} has output
error at most $\rho=[C L(a_0+1)]^{-1}$. It suffices that
\[
 \log(1/\Delta)\le
 C[d+\log(L(a_0+2))+\sqrt{L(a_0+2)}]
\]
for an appropriate sufficiently small tolerance. The required width
in~\eqref{ac-dual-width} is therefore at most $2^{Cd}R_0^C$.
The repaired output keeps the selected signs; multiplying its readout
by $(1-2\rho)^{-1}$ restores margin $1/2$ and adds $O(1/L)$
action. This realizes a function in $E$ at action at most
$a_0+K_1/L$, contradicting $K_0>K_1$.

Hence $\mu_N(E)<e^{-N(a_0+1/L)}$. Division by the sample-consistent
denominator bounds posterior bad mass by $e^{-N/L}\le1/32$ under
the stated width threshold. Absorb $K_0/L$ into the logarithmic
excess. Theorem~\ref{ac-circuits} gives the claimed circuit bounds.
If the rate infimum is not attained, choose a path within $1/H$;
this changes the exponent $\sqrt{H\epsilon}$ only by a constant.
Every bound is simultaneous over all samples of the specified size,
and the class itself is sample independent.
\end{proof}
